% Current local-window proof. Compatible with the general comparison in the following section.
\section{三つのstabilizationによる局所残余圏の表示}
\label{tw:sec:window}

本節は、主定理で用いる局所比較(L)を、代数的入力から導く。
形式的MF同値(CF)、相対対角のperfect性、特殊ファイバーからの生成性の持上げは使わない。
完備化はHom複体の有限性を証明するためだけに用い、
関手自体は代数的なsection--windowと余核のdg比較から作る。
したがって、両立するArtin同値の列を逆極限にして関手を得たとは主張しない。

$(R,\mathfrak m,k)$を有限Krull次元の正則noetherian局所$\mathbb C$-代数、
$k=\mathbb C$とし、$\widehat R$をその完備化とする。
$A=\bigoplus_{j\ge0}A_j$は中央$R$-代数、$A_0=R$、
各$A_j$は有限自由で、$A\otimes_R k=k[x_1,\ldots,x_N]$、$N\ge5$とする。
$\Omega\in Z(A)_3$の剰余$F$は滑らかな射影三次超曲面を定めるものとし、
$B=A/(\Omega)$、$a=N-3$、$i=1-a$と置く。
完備化後も$A_{\widehat R}=\bigoplus_j(A_j\otimes_R\widehat R)$という
\emph{次数別}の代数を使い、内部次数についての無限級数は導入しない。

\subsection{完備Hom複体の縮約と代数的基底への降下}

\begin{lemma}[次数別完備な複体の有限縮約]\label{tw:lem:contraction}
$S$を剰余体$\mathbb C$の完備noetherian局所$\mathbb C$-代数とし、
$C^\bullet$を各項が有限自由$S$-加群である複体とする。
$C^\bullet\otimes_S\mathbb C$の全cohomologyの次元の和が有限なら、
$C^\bullet$は有界有限自由複体と$S$-線形にhomotopy同値である。
複体の全長が有限であることは仮定しない。
\end{lemma}
\begin{proof}
$C_0=C\otimes_S\mathbb C$のcohomologyを$H$とする。
体上の分裂から、次数を保つ$i_0:H\to C_0$、$p_0:C_0\to H$と
次数$-1$の$s_0$を
\begin{equation}\label{tw:eq:retract}
 p_0i_0=1,\qquad d_0s_0+s_0d_0=1-i_0p_0,
 \qquad s_0^2=p_0s_0=s_0i_0=0
\end{equation}
となるように選べる。各次数の基底を選び、これらを$S$へ延長する。
元の微分は$d=d_0+\delta$、$\delta\in\mathfrak m_S\operatorname{Hom}^1(C,C)$と書ける。
各固定次数は有限自由なので、次の逆行列はその次数内で収束する：
\begin{equation}\label{tw:eq:perturbation}
 \begin{aligned}
 J&=(1+s_0\delta)^{-1},& L&=(1+\delta s_0)^{-1},\\
 i'&=Ji_0,&p'&=p_0L,\\
 s'&=s_0L,&d_H&=p_0\delta Ji_0.
 \end{aligned}
\end{equation}
これは標準のhomological perturbation公式である
\cite[Theorem~2.3]{TWCrainic}。
同文献のhomotopyの符号は本稿の$-s_0$に対応する。
$d^2=0$と式\eqref{tw:eq:retract}を用いると
$d_H^2=0$、$di'=i'd_H$、$p'd=d_Hp'$、$p'i'=1$、
$ds'+s'd=1-i'p'$となる。
各積は固定次数での有限行列の積であり、逆行列の級数だけが$\mathfrak m_S$進級数である。
$H$は有界有限次元なので$(H\otimes_{\mathbb C}S,d_H)$が所要の有限複体である。
\end{proof}

\begin{lemma}[有限rank MFのproper Homと底変更]\label{tw:lem:proper-hom}
上の$R,A,\Omega$に対し、有限自由graded MFの任意の二対象$P,Q$の
Hom複体は$R$-perfectである。その通常の複体は有界有限自由複体とhomotopy同値であり、
$R\to k,\widehat R,\operatorname{Frac}(R)$について
\begin{equation}\label{tw:eq:hom-bc}
 \operatorname{Hom}_{\mathrm{MF}}(P,Q)\otimes_R^{\mathbf L}S
 \simeq\operatorname{Hom}_{\mathrm{MF}_S}(P_S,Q_S)
\end{equation}
となる。ここで右辺は有限rank行列の通常の底変更から作るHomである。
\end{lemma}
\begin{proof}
周期変数の二重次数を$(2,-3)$とする。
固定cohomological degreeのHom行列成分では周期変数の冪と内部次数が固定される。
従ってHom複体$C$の各項は有限個の$A_j$の直和、すなわち有限自由$R$-加群であり、
各項の形成と通常の底変更は可換である。
剰余のHomは滑らかな可換cubicのgraded MFのHomなので、
古典Orlov比較により有界有限次元cohomologyを持つ。
補題\ref{tw:lem:contraction}を$C\otimes_R\widehat R$へ適用すると、
完備化後の複体は有界有限自由複体へ縮約する。

$R$はnoetherianなので$H^q(C)$は各$q$で有限生成である。
平坦な完備化はcohomologyと可換であり、忠実平坦性から
$H^q(C)=0$は$q$の有限区間の外で成り立つ。
正則$R$の大域次元は有限なので$C$は$R$-perfectである。
さらに有界有限自由複体$D$から$C$への準同型を選ぶと、そのconeは
各項が射影である完全複体になる。
$e=\operatorname{gldim}R$とすれば、完全複体のcycleに対して
$\operatorname{Ext}^1_R(Z^q,M)=\operatorname{Ext}^{e+1}_R(Z^{q+e},M)=0$
というdimension shiftingが適用できる。よって全cycleは射影で、このconeは可縮である。
従って$C$自体が$D$とhomotopy同値、特にK-flatである。
項ごとの底変更の同定が式\eqref{tw:eq:hom-bc}を与える。
有限coneと直和因子にも同じ結論が延長される。
\end{proof}

\subsection{stabilizationとsection--windowの実際のenhancement}

\begin{lemma}[正則局所基底上のdg stabilization]\label{tw:lem:stabilization}
$A,B$は左右noetherian、$A$の大域次元は有限、$\Omega$は正則である。
余核による比較は、内部shiftと中心元に由来する$(3)\simeq[2]$を保つ
\begin{equation}\label{tw:eq:stabilization}
 \mathrm{MF}^{\mathbb Z}(A,\Omega)
 \simeq D^b_{\mathrm{sg,dg}}(\operatorname{gr}B)^{\mathrm{idem}}
\end{equation}
というdg Morita同値を与える。
本節の有限windowに現れる対象では、$k$および$\widehat R$への底変更と整合する。
\end{lemma}
\begin{proof}
noetherian性、相対Koszul分解、$\operatorname{gldim}A\le N+\dim R$、
$\Omega$の正則性、$B_j$の自由性と有限左右injective dimensionは、
補題\ref{lp:lem:noetherian}、\ref{lp:lem:resolution}と
命題\ref{lp:prop:koszul}、\ref{lp:prop:gorenstein}による。
これらの次数別の証明は基底の完備性や有限型性を要求せず、$\widehat R$でも通る。

Chenの余核定理\cite[Theorem~7.2]{AMChen}は三角圏の定理である。
そのままdg同値を引用するのではなく、以下のenhancementを用いる。
有限自由graded MFの成分別split exact圏を$\mathcal F$、
有限生成graded Gorenstein-projective $B$-加群のexact圏を$\mathcal G$とする。
$A$の大域次元が有限なので、$\mathcal G$の全対象は有限$A$-射影次元を持つ。
ChenのCorollary~2.11により、非零$M\in\mathcal G$の$A$-射影次元は一である
（零対象も含めれば高々一）。従って有限生成元を斉次に選ぶと
\begin{equation}\label{tw:eq:graded-presentation}
 0\longrightarrow P_0\xrightarrow{d_0}P_1\longrightarrow M\longrightarrow0
\end{equation}
という有限graded projective $A$-加群による表示を得る。
実際、有限graded freeな$P_1$からの全射の核は$A$上射影であり、
斉次写像の持上げの次数零成分を取ることでgradedにも射影である。
各有限graded projective $A$-加群は、$A_0=R$が局所であることから
有限個の内部shiftを持つ自由加群である。
$\Omega M=0$より、次数零の乗法写像$\Omega_{P_1}:P_1\to P_1(3)$は
$d_0(3)$を経由し、一意な$d_1:P_1\to P_0(3)$を定める。
従って
\begin{equation}\label{tw:eq:graded-factorization}
 d_0(3)d_1=\Omega_{P_1},\qquad d_1d_0=\Omega_{P_0}.
\end{equation}
第二式は第一式と$d_0(3)$の単射性から従う。
これが本稿の内部shift規約における有限自由graded MFで、その余核は$M$である。
この構成はChenのTheorem~7.2とRemark~7.3(2)の本質的全射性を斉次に実現する。
余核はexact関手$\mathcal F\to\mathcal G$であり、projective-injectiveを保つ。
射とhomotopyの持上げも次数零成分を取れるので、同定理の完全忠実性もgraded版で成立する。
右加群については$A^{\mathrm{op}}$へ同じ結果を適用する。
従って、その誘導するstable categoryの関手は同値である。

exact圏の有界導来dg圏をprojective-injectiveの有界複体で割ることで、
実際のdg関手のzigzagを得る。
左側とMFの比較は有限totalizationである。
成分別splitな完全複体に横方向のcontraction $s$を取ると、
縦MF微分$\partial$に対する補正$s(1+\partial s)^{-1}$は横方向の有限性により有限和である。
中心曲率$w$の残差は$sws=ws^2=0$なので、totalizationは可縮である。
右側と特異圏の比較は$\mathcal G$の包含から誘導される。
完全射影分解がstable Homを計算し、有限injective dimensionによる高いsyzygyが
本質的全射性を与える。各zigzagの$H^0$は上記の三角同値である。
全てのshiftにも同じ議論を適用することでHomの全cohomologyの同型となり、
式\eqref{tw:eq:stabilization}を得る。

指定周期の比較は、余核の二拡張で確認できる。
式\eqref{tw:eq:graded-factorization}のMFの余核を$M$とし、
$\overline P_j=P_j\otimes_A B$と書く。その周期分解の連続する二段は
\begin{equation}\label{tw:eq:period-extension}
 0\longrightarrow M(-3)\xrightarrow{\overline d_1(-3)}
 \overline P_0\xrightarrow{\overline d_0}\overline P_1
 \longrightarrow M\longrightarrow0
\end{equation}
である。左端の写像は$d_1(-3)$が余核を経由して誘導する写像を表す。
二つの微分を$A$へ持ち上げた合成は、式\eqref{tw:eq:graded-factorization}により
指定した$\Omega$の乗法そのものである。
従って、$\Omega\bmod(\Omega)^2$を用いて
$(\Omega)/(\Omega)^2=B(-3)$と同定した導来自己交叉の連結射
$\eta_M:M\to M(-3)[2]$は、この二拡張のYoneda類である。
ここでは命題\ref{lp:prop:sing-period}と同じconeの符号規約を用いる。
したがって余核は、MFの二段周期を
\begin{equation}\label{tw:eq:specified-period}
 \tau_M^{\rm sg}:=\eta_M(3):M(3)\xrightarrow{\sim}M[2]
\end{equation}
へ移す。有限totalizationは内部shiftと複体の二重shiftに可換であり、
同じ二拡張をtotalizeするので、この比較を保つ。
$\mathcal G$から特異圏への包含も二拡張をそのまま送る。
これで全てのzigzagに沿う周期の規格化が決まり、任意のスカラー倍を挿入しない。

底変更には余核構成を使う。有限MFの二つの行列の積は$\Omega$であり、
剰余では非零多項式$F$の乗法になるため、各行列は剰余上単射である。
固定内部次数で単元小行列式を選ぶと、余核は有限自由$R$-加群となる。
従って余核の周期分解、有限totalization、windowで使う次数別自由な有限長加群は
通常の剰余底変更と整合する。完備化はさらに平坦である。
Morita逆であるstabilizationの比較は、この余核の自然な比較の逆として選ぶ。
任意のcoherent対象の任意の非正則底変更まで主張してはいない。
\end{proof}

\subsection{比較図式と関手の定義域}\label{tw:subsec:comparison}
上のexact圏$\mathcal F,\mathcal G$のprojective-injective部分圏を
$\mathcal P_{\mathcal F},\mathcal P_{\mathcal G}$とし、
\[
 \mathfrak F=
 \bigl(D^b_{\rm dg}(\mathcal F)/K^b_{\rm dg}(\mathcal P_{\mathcal F})\bigr)^{\rm idem},
 \qquad
 \mathfrak G=
 \bigl(D^b_{\rm dg}(\mathcal G)/K^b_{\rm dg}(\mathcal P_{\mathcal G})\bigr)^{\rm idem}
\]
と置く。ここで$D^b_{\rm dg}$はexact圏の有界導来dg圏、斜線はdg商を表す。
さらに$\mathcal D_B=D^b_{\rm dg}(\operatorname{gr}B)$、
$\mathcal S_B=(\mathcal D_B/\operatorname{Perf}_{\rm dg}(B))^{\rm idem}$、
$\mathcal M_R=\mathrm{MF}^{\mathbb Z}(A,\Omega)$と略記する。
$\mathcal M_R$にも有限coneと冪等完備化を施す。

derived sectionの切断を$\sigma_i=(R\omega(-))_{\ge i}$、
例外成分$\mathcal E_B$を除くmutationを$p$とすると、
stabilizationと次のwindow定理の関手は一つの図式に収まる：
\begin{equation}\label{tw:eq:comparison-diagram}
 \begin{array}{c@{\ }c@{\ }c@{\ }c@{\ }c@{\ }c@{\ }c}
 \mathcal C_B&\xrightarrow{\sigma_i}&\mathcal D_B&\xrightarrow{q}&
 \mathcal S_B&\xleftarrow[\sim]{\overline\iota}&\mathfrak G\\[2pt]
 {\scriptstyle p}\!\downarrow&&&&&&\Vert\\[2pt]
 \mathcal K_B^{\rm lin}&\xrightarrow{\Psi_R}&\mathcal M_R&
 \xleftarrow[\sim]{\operatorname{Tot}}&\mathfrak F&
 \xrightarrow[\sim]{\overline c}&\mathfrak G .
 \end{array}
\end{equation}
ここで$q$はperfect複体によるdg商、$\overline c$は余核exact関手から、
$\overline\iota$は$\mathcal G\hookrightarrow\operatorname{gr}B$から誘導する。
$\operatorname{Tot}$は有界複体の有限totalizationである。
上段の$\sigma_i$の定義域は線束生成部分$\mathcal C_B$であり、
任意のcoherent対象に対する底変更をこの図式で主張してはいない。
記号$\sim$を付けた矢印は補題\ref{tw:lem:stabilization}のdg Morita同値である。
図式の可換性はMorita局所化における自然な比較
\begin{equation}\label{tw:eq:comparison-composite}
 q\sigma_i\simeq
 \overline\iota\,\overline c\,\operatorname{Tot}^{-1}\Psi_Rp
\end{equation}
を意味し、$\Psi_R$の像と完全忠実性は以下で証明する。
式\eqref{tw:eq:period-extension}は
$\operatorname{Tot},\overline c,\overline\iota$の各矢印に沿う指定周期を与える。
$\mathcal C_B$と$\mathcal D_B$自体に三次周期性があるとは述べない。
$q\sigma_i$を残余圏へ制限した後のrotationとの比較は
式\eqref{tw:eq:rotation-period-square}で明示する。

\begin{theorem}[三対象window定理：局所比較(L)の確立]\label{tw:thm:window}
本節の代数的入力だけを仮定する。$K_R=\operatorname{stab}(R)$、
$T_R=K_R\oplus K_R(1)\oplus K_R(2)$と置くと、
\begin{equation}\label{tw:eq:main}
 \mathcal K_{B_R}^{\mathrm{lin}}
 \simeq\operatorname{thick}\{K_R,K_R(1),K_R(2)\}
 \simeq\operatorname{Perf}(\Gamma_R),\qquad
 \Gamma_R=R\operatorname{End}(T_R)\in\operatorname{Perf}(R)
\end{equation}
である。最初の同値は
\begin{equation}\label{tw:eq:window}
 \Psi_R=\operatorname{stab}\circ\bigl((R\omega -)_{\ge i}\bigr),
 \qquad i=1-a
\end{equation}
の制限で与えられる。
これは$k,\widehat R,\operatorname{Frac}(R)$への底変更と整合し、
$L_{\mathcal O}\circ(1)$をMFの内部shiftへ移す。
特にこのrotationは自己同値であり、中心元に由来する三次周期同型を持つ。
\end{theorem}
\begin{proof}
$E_R=\bigoplus_{j=0}^{N-1}\mathcal O(j)$は線束生成部分を生成し、
その全対象対のHomは$R$-perfectである。
例外列$\mathcal E=\langle\mathcal O,\ldots,\mathcal O(a-1)\rangle$を除くmutationを$p$とすると、
$G_R=p(E_R)$が$\mathcal K^{\mathrm{lin}}$を生成する
（命題\ref{lp:prop:residual}）。
必要なsection、有限cohomological dimension、$\chi$の入力は
Burke--Stevenson\cite[Proposition~4.18, Lemmas~6.1--6.2]{LP-BS}による。

\emph{関手の定義。}
線束Homの公式から、$0\le j\le N-1$では
$(R\omega\mathcal O(j))_{\ge i}=B(j)_{\ge i}$である。
$T_j=B(j)_{<i}$とすると、$T_j$の各次数は有限自由で、
$R$のshiftによる有限filtrationを持つ。
$R$の相対Koszul分解と$B$の長さ一の$A$-分解により、
$B(j)_{\ge i}$は$A$-perfectである。
従ってderived section、内部切断、dg商、補題\ref{tw:lem:stabilization}のMorita逆を合成して
式\eqref{tw:eq:window}が実際のdg quasi-functorとして定義される。
$0\le j<a$では$T_j=0$なので$\Psi_R(\mathcal O(j))=0$であり、
$\Psi_Rp\simeq\Psi_R$である。

\emph{完全忠実性。}
この関手の剰余は、古典的なsection--windowを特異圏へ送る関手である。
Burke--Stevenson\cite[Theorem~6.4とその証明]{LP-BS}において$i=1-a$を取ると、
そのkernelは$\mathcal E_k$であり、Verdier商を特異圏と同定する。
本稿の右直交と同文献の半直交分解の向きは異なり得るが、
双方が同じVerdier商を表示するため、$\mathcal E_k^\perp$への制限は同値である。

sourceの線束Homの底変更は命題\ref{lp:prop:basechange}による。
sectionの内部切断は生成線束上で上の加群の表示になり、
stabilizationの底変更は補題\ref{tw:lem:stabilization}による。
有限cone、直和因子、evaluationを通じ、$\Psi_R$のHom上の\emph{自然な写像}は
剰余で古典windowのHom写像へ移る。
targetのHomは補題\ref{tw:lem:proper-hom}により$R$-perfectかつ剰余底変更と可換である。
従って任意の二対象についてHom比較のconeはperfectで剰余が零となり、
局所環上の導来Nakayamaにより零である。これがenhancedな完全忠実性を与える。
単にHomの階数だけを比較しているのではない。

\emph{像の同定。}
残る$j=a,a+1,a+2$について、$B(j)$は特異圏で零なので
\begin{equation}\label{tw:eq:triangular}
 \Psi_R(\mathcal O(j))\simeq\operatorname{stab}(T_j)[-1],\qquad
 \operatorname{gr}(T_j)=\bigoplus_{r=0}^{j-a}B_r\otimes_RR(j-r).
\end{equation}
$B_0=R$である。まず$T_a=R(a)$、次に$T_{a+1}$の高次数部分は
$R(a)$の有限直和、最後に$T_{a+2}$の高次数部分は
$R(a),R(a+1)$の有限extensionである。
その商はそれぞれ$R(a+1),R(a+2)$なので、三角的に
\begin{equation}\label{tw:eq:three-thick}
 \operatorname{thick}(\Psi_R(G_R))
 =\operatorname{thick}\{K_R(a),K_R(a+1),K_R(a+2)\}
 =\operatorname{thick}(T_R).
\end{equation}
最後はMFの$(3)=[2]$による。
完全忠実なexact関手の冪等完備sourceの像はthickなので、これが式\eqref{tw:eq:main}の同値となる。
$\Gamma_R$のproper性は補題\ref{tw:lem:proper-hom}からも従う。
生成性の議論は有限長加群の実際のfiltrationによるものであり、
特殊ファイバーの生成性だけを持ち上げたものではない。

\emph{rotationと三次周期性。}
線束Homの消滅とevaluation三角により、$L_{\mathcal O}(M(1))$は再び残余圏に属する
（命題\ref{lp:prop:rotation-candidate}）。
$M\in\mathcal K^{\mathrm{lin}}$なら
$(R\omega M)_i=R\operatorname{Hom}(\mathcal O(a-1),M)=0$である。
切断の自然な三角により
\begin{equation}\label{tw:eq:rotation-cut}
 ((R\omega M)(1))_{\ge i}
 =(R\omega M)_{\ge i+1}(1)
 \simeq(R\omega M)_{\ge i}(1).
\end{equation}
$\Psi_R(\mathcal O)=0$とevaluationの三角を使うと
$\Psi_R(L_{\mathcal O}M(1))\simeq\Psi_R(M)(1)$を得る。
式\eqref{tw:eq:three-thick}の像は$(1)$と$(-1)$の両方で安定なので、
rotationは自己同値となる。
周期同型はMFで中心元$\Omega$から定まる$(3)=[2]$をこの同値で移したものとする。
より正確には、上で得た自然比較を
$\beta_M:\Psi_R(\rho_BM)\xrightarrow{\sim}\Psi_R(M)(1)$とし、
$\beta_M^{(3)}=\beta_M(2)\circ\beta_{\rho_BM}(1)\circ\beta_{\rho_B^2M}$と置く。
MFの指定周期を$\tau^{\rm MF}$と書くと、$\Psi_R$の完全忠実性により
次の正方形を可換にする$\tau^{\rm rot}$が一意に定まる：
\begin{equation}\label{tw:eq:rotation-period-square}
 \begin{array}{ccc}
 \Psi_R(\rho_B^3M)&\xrightarrow{\Psi_R(\tau_M^{\rm rot})}&\Psi_R(M[2])\\[2pt]
 {\scriptstyle\beta_M^{(3)}}\!\downarrow&&\downarrow\!{\scriptstyle\simeq}\\[2pt]
 \Psi_R(M)(3)&\xrightarrow{\tau_{\Psi_R(M)}^{\rm MF}}&\Psi_R(M)[2].
 \end{array}
\end{equation}
図式\eqref{tw:eq:comparison-diagram}の右側では、下辺の周期は
$\tau^{\rm sg}=\eta(3)$へ移る。
従ってこの正方形は、切断とevaluationが与えるrotation比較と、
中心元の二拡張が与える周期比較とを同時に固定している。
各矢印は余核・切断・evaluationの自然な写像から作ったので、
以下の底変更とも整合する。別に選んだ周期のスカラー倍を保存するとは主張しない。

完備化・分数体への底変更は平坦であり、有限MF行列のHomと積に項ごとに可換する。
sourceも有限線束生成とmutationから作った模型の底変更に一致する。
剰余底変更は上で確認した。これで必要な三対象とrotationの比較をすべて得る。
\end{proof}
